% CREACION DEL DOCUMENTO
\documentclass[
	spanish,
	letterpaper, oneside
]{article}

% INFORMACION DEL DOCUMENTO
\def\documenttitle {Linea de tiempo: autores sobre la etica en Mexico}
\def\documentsubtitle {Principales autores sobre la etica en Mexico}
\def\documentsubject {Actividad 3 - Etica y Moral juridica}

\def\documentauthor {Martin Jonathan de la Cruz}
\def\coursename {Etica y Moral juridica}
\def\coursecode {030}

\def\universityname {Universidad Abierta y a Distancia de Mexico}
\def\universityfaculty {Licenciatura en Derecho}
\def\universitydepartment {Etica y Moral juridica}
\def\universitydepartmentimage {departamentos/UnADM}
\def\universitydepartmentimagecfg {height=1.57cm}
\def\universitylocation {Roma Norte, Ciudad de Mexico}

% INTEGRANTES, PROFESORES Y FECHAS
\def\authortable {
	\begin{tabular}{ll}
		\textbf{Alumno:}
		& \begin{tabular}[t]{l}
			 Martin Jonathan de la Cruz \\
		\end{tabular} \\
		Matricula: & ES2611202040 \\

		\textit{Profesor:}
		& \begin{tabular}[t]{l}
			Emma Liliana \\ Padilla Cano
		\end{tabular} \\
		& \\
		\multicolumn{2}{l}{Fecha de realizacion: \today} \\
		\multicolumn{2}{l}{\universitylocation}
	\end{tabular}
}

% IMPORTACION DEL TEMPLATE
\input{template}

% FORMATO SOLICITADO
\usepackage{helvet}
\renewcommand{\familydefault}{\sfdefault}

% PRODUCTO VISUAL
\usepackage{tikz}
\usetikzlibrary{arrows.meta,positioning,calc,shadows.blur}

% CITAS EN FORMATO AUTOR-ANIO
\setcitestyle{authoryear,open={(},close={)}}

% ==============================================================================
% MARCA DE AGUA EN PORTADA
% - Se aplica solo a la portada porque usa \AddToShipoutPictureBG*.
% - Para apagarla en alguna actividad: \def\coverwatermarkenabled {false}
% - Usar opacidad baja para que no compita con titulo, datos ni logotipo.
% ==============================================================================
\def\coverwatermarkenabled {true}
\def\coverwatermarkimage {img/departamentos/UnADM.pdf}
\def\coverwatermarkopacity {0.16}
\def\coverwatermarkwidth {0.72\paperwidth}
\def\coverwatermarkxshift {0cm}
\def\coverwatermarkyshift {-0.35cm}

\newcommand{\insertcoverwatermark}{%
	\ifthenelse{\equal{\coverwatermarkenabled}{true}}{%
		\AddToShipoutPictureBG*{%
			\begin{tikzpicture}[remember picture,overlay]
				\node[
					opacity=\coverwatermarkopacity,
					inner sep=0pt,
					xshift=\coverwatermarkxshift,
					yshift=\coverwatermarkyshift
				] at (current page.center) {%
					\includegraphics[width=\coverwatermarkwidth]{\coverwatermarkimage}%
				};
			\end{tikzpicture}%
		}%
	}{}%
}

% Tarjeta compacta para la linea de tiempo: retrato + autor + fecha + aporte.
\newcommand{\timelineauthor}[4]{%
	\begingroup
	\setlength{\tabcolsep}{2pt}
	\begin{tabular}{@{}m{0.85cm}p{2.95cm}@{}}
		\includegraphics[width=0.8cm,height=0.9cm,keepaspectratio]{#1}
		& {\textbf{#2}\par\textbf{#3}}
	\end{tabular}\\[-1pt]
	\endgroup
	#4%
}

\begin{document}

% PORTADA
\insertcoverwatermark
\templatePortrait

% CONFIGURACION DE PAGINA Y ENCABEZADOS
\templatePagecfg
\onehalfspacing

% RESUMEN
\begin{abstractd}
	\textbf{Objetivo:} Identificar y analizar las aportaciones de autores relevantes para la reflexion etica en Mexico, desde el pensamiento filosofico-educativo del siglo XIX y XX hasta enfoques contemporaneos vinculados con la etica profesional, politica y bioetica.

	\textbf{Producto elaborado:} Se presenta una linea de tiempo en TikZ titulada ``Principales autores sobre la etica en Mexico'', organizada con fechas, autores, campos de aplicacion y aportaciones principales.

	\textbf{Enfoque de analisis:} La seleccion no se entiende como una lista de nombres, sino como un recorrido de problemas: educar moralmente, formar ciudadania, pensar la identidad, criticar la injusticia, orientar la responsabilidad profesional y proteger la dignidad humana.
\end{abstractd}

% TABLA DE CONTENIDOS
\templateIndex

% CONFIGURACIONES FINALES
\templateFinalcfg

\section{Introduccion}

La etica en Mexico no se ha desarrollado como una reflexion aislada de la vida publica. Al contrario, aparece ligada a problemas concretos: la educacion del ciudadano, la construccion de una identidad nacional, la critica de la desigualdad, la responsabilidad profesional y la defensa de la dignidad humana. Por eso, estudiar a sus principales autores no consiste solamente en ubicar fechas, sino en reconocer como cada momento historico produjo una pregunta moral distinta.

La \textbf{etica}, entendida como reflexion racional sobre la \textbf{conducta moral}, permite distinguir entre costumbre, norma, valor y responsabilidad. Martinez Huerta explica que la etica no se reduce a repetir la moral vivida, sino que busca razones para comprender y valorar la accion humana \citep{huertaEticaConClasicos2000}. En sentido semejante, Ronquillo Armas vincula la etica con la conducta humana, la moralidad, los valores y la \textbf{responsabilidad profesional} \citep{ronquilloarmasEticaGeneralProfesional2018}. Por tanto, sostengo que estas nociones permiten leer el caso mexicano no como una sucesion de autores, sino como una reflexion historicamente situada sobre la formacion moral de la persona y la responsabilidad social de sus actos.

% La postura que guia este trabajo es que la reflexion etica mexicana ha pasado de una preocupacion educativa y civica hacia una reflexion mas amplia sobre justicia, praxis, comunidad, profesion y bioetica. Ese cambio importa para la formacion juridica porque el Derecho no solo aplica reglas: tambien interpreta valores, protege personas y decide frente a conflictos reales.

% \section{Criterio de seleccion y tecnica didactica}

\section{Línea de tiempo de los principales autores sobre la ética en México}

% La linea de tiempo se elaboro con la tecnica didactica de organizacion cronologica, util para ordenar procesos historicos, reconocer cambios conceptuales y relacionar autores con contextos. Esta tecnica permite que la informacion no aparezca como datos sueltos, sino como secuencia razonada: periodo, autor, aportacion, ambito de aplicacion y relevancia etico-juridica \citep{lopezmartinezTecnicasDidacticas2023}.

% Para cumplir con la actividad, 
Se eligieron ocho autores cuyo pensamiento permite observar una evolucion amplia de la etica en Mexico: Gabino Barreda, Justo Sierra, Antonio Caso, Jose Vasconcelos, Samuel Ramos, Adolfo Sanchez Vazquez, Luis Villoro y Juliana Gonzalez Valenzuela. La seleccion combina pensamiento educativo, filosofia moral, identidad cultural, etica de la praxis, etica politica y bioetica. Aunque no todos escribieron exclusivamente tratados de etica juridica, sus aportaciones ayudan a comprender el trasfondo moral desde el cual se forman criterios de ciudadania, justicia, dignidad y responsabilidad.

Las fechas de la linea de tiempo toman como referencia obras, discursos o publicaciones representativas de cada autor: la \textit{Oracion civica} de Barreda, el discurso universitario de Sierra, la obra humanista de Caso, la reflexion cultural de Vasconcelos y Ramos, la etica de la praxis de Sanchez Vazquez, la filosofia del ethos de Gonzalez Valenzuela y la etica politica de Villoro \citep{barredaOracionCivica1867,sierraUniversidadNacional1910,casoExistencia1916,vasconcelosRazaCosmica1925,ramosPerfil1934,sanchezVazquezEtica1969,gonzalezEthosDestino1996,villoroPoderValor1997}.

\clearpage
\pagestyle{empty}
\begin{landscape}
\begin{figure}[H]
\centering
\resizebox{0.98\linewidth}{!}{%
\begin{tikzpicture}[
	timeline/.style={-{Stealth[length=3mm]}, very thick, black!70},
	tmark/.style={circle, draw=black, fill=black, inner sep=2.2pt},
	cardtop/.style={
		rectangle,
		rounded corners=2pt,
		draw=black!65,
		fill=gray!8,
		text width=4.35cm,
		align=left,
		inner sep=6pt,
		font=\scriptsize
	},
	cardbottom/.style={
		rectangle,
		rounded corners=2pt,
		draw=black!65,
		fill=gray!3,
		text width=4.35cm,
		align=left,
		inner sep=6pt,
		font=\scriptsize
	},
	year/.style={font=\small\bfseries, align=center}
]
\draw[timeline] (-0.5,0) -- (30,0);

\node[tmark] at (0,0) {};
\node[year] at (0,0.45) {1867};
\draw[black!45, thick] (0,0) -- (0,0.95);
\node[cardtop] at (0,3.1) {\timelineauthor{img/autores_etica/gabino_barreda.jpg}{Gabino Barreda}{1867}{Introduce el positivismo como proyecto educativo y civico. Su aporte consistio en vincular orden, razon y progreso con la formacion moral del ciudadano despues de la Reforma. Ambito: educacion publica y moral civica.}};

\node[tmark] at (4.2,0) {};
\node[year] at (4.2,0.45) {1910};
\draw[black!45, thick] (4.2,0) -- (4.2,-0.95);
\node[cardbottom] at (4.2,-3.25) {\timelineauthor{img/autores_etica/justo_sierra.jpg}{Justo Sierra}{1910}{Defendio una universidad orientada a formar conciencia nacional, libertad intelectual y servicio publico. Su etica educativa conecta cultura, responsabilidad social y formacion profesional. Ambito: educacion superior y ciudadania.}};

\node[tmark] at (8.4,0) {};
\node[year] at (8.4,0.45) {1916};
\draw[black!45, thick] (8.4,0) -- (8.4,0.95);
\node[cardtop] at (8.4,3.1) {\timelineauthor{img/autores_etica/antonio_caso.jpg}{Antonio Caso}{1916}{Critico el reduccionismo positivista y propuso una vision humanista de la persona. Su idea de desinteres y caridad desplaza el calculo utilitario hacia la responsabilidad por el otro. Ambito: filosofia moral y humanismo.}};

\node[tmark] at (12.6,0) {};
\node[year] at (12.6,0.45) {1925};
\draw[black!45, thick] (12.6,0) -- (12.6,-0.95);
\node[cardbottom] at (12.6,-3.25) {\timelineauthor{img/autores_etica/jose_vasconcelos.jpg}{Jose Vasconcelos}{1925}{Relaciono educacion, cultura e identidad nacional con un proyecto moral de integracion. Su pensamiento impulsa una etica humanista de la cultura, aunque exige lectura critica por sus tensiones historicas. Ambito: educacion y cultura.}};

\node[tmark] at (16.8,0) {};
\node[year] at (16.8,0.45) {1934};
\draw[black!45, thick] (16.8,0) -- (16.8,0.95);
\node[cardtop] at (16.8,3.1) {\timelineauthor{img/autores_etica/samuel_ramos.jpg}{Samuel Ramos}{1934}{Analizo el caracter mexicano y sus conflictos de autovaloracion. Su aporte etico consiste en mirar la moral social desde la identidad, la autenticidad y la necesidad de superar complejos culturales. Ambito: cultura e identidad.}};

\node[tmark] at (21.0,0) {};
\node[year] at (21.0,0.45) {1969};
\draw[black!45, thick] (21.0,0) -- (21.0,-0.95);
\node[cardbottom] at (21.0,-3.25) {\timelineauthor{img/autores_etica/adolfo_sanchez_vazquez.jpg}{Adolfo Sanchez Vazquez}{1969}{Entendio la etica como teoria del comportamiento moral en sociedad y como praxis transformadora. Su enfoque permite vincular libertad, responsabilidad y critica de condiciones injustas. Ambito: etica social y politica.}};

\node[tmark] at (25.2,0) {};
\node[year] at (25.2,0.45) {1996};
\draw[black!45, thick] (25.2,0) -- (25.2,0.95);
\node[cardtop] at (25.2,3.1) {\timelineauthor{img/autores_etica/juliana_gonzalez.jpg}{Juliana Gonzalez}{1996}{Desarrollo una reflexion sobre ethos, libertad, dignidad y responsabilidad humana. Su pensamiento abre camino para discutir bioetica y cuidado de la vida desde una filosofia mexicana contemporanea. Ambito: bioetica y humanismo.}};

\node[tmark] at (29.4,0) {};
\node[year] at (29.4,0.45) {1997};
\draw[black!45, thick] (29.4,0) -- (29.4,-0.95);
\node[cardbottom] at (29.4,-3.25) {\timelineauthor{img/autores_etica/luis_villoro.jpg}{Luis Villoro}{1997}{Planteo una etica politica centrada en valor, poder, justicia y comunidad. Su obra ayuda a pensar la legitimidad de las instituciones y la critica de la dominacion. Ambito: filosofia politica y justicia.}};
\end{tikzpicture}}
\caption{Principales autores sobre la etica en Mexico}
\label{fig:linea-etica-mexico}
\end{figure}
\end{landscape}
\pagestyle{fancy}
\clearpage

\section{Evolucion de la educacion moral y la ciudadania}

\subsection{Educación cívica y responsabilidad pública}

La línea de tiempo muestra una primera etapa marcada por la educación moral y la organización cívica. Puede entenderse la \textbf{educación moral} como la formación de criterios, hábitos y responsabilidades que orientan la conducta en la vida común. En Barreda y Sierra, la ética aparece ligada a la formación del ciudadano: educar no significa solo transmitir conocimientos, sino desarrollar disciplina intelectual y sentido de responsabilidad pública, como se advierte en sus obras y discursos representativos \citep{barredaOracionCivica1867,sierraUniversidadNacional1910}. \textbf{Desde mi perspectiva}, esta etapa revela algo importante para el Derecho: una sociedad no se sostiene únicamente con códigos, sino con personas capaces de reconocer razones, límites y deberes compartidos.

% Se enumeran los tres desplazamientos porque cada uno responde a una pregunta
% distinta sobre el vínculo entre ética y derecho.
\begin{enumerate}
  \item \textbf{De la moral privada a la responsabilidad pública}: la ética deja de ser asunto de conciencia individual y pasa a fundar deberes cívicos exigibles \citep{barredaOracionCivica1867}.
  \item \textbf{De la instrucción a la formación del criterio}: el fin educativo se desplaza de transmitir contenidos a desarrollar juicio propio \citep{sierraUniversidadNacional1910}.
  \item \textbf{Del deber abstracto a la práctica institucional}: los principios éticos requieren instituciones que los vuelvan operativos, y esa es la aportación específica del Derecho.
\end{enumerate}

\subsection{Humanismo, cultura e identidad}

En la etapa humanista, el \textbf{humanismo ético} puede entenderse como una perspectiva que valora la acción a partir de la persona, su dignidad y los fines que orientan su vida en común. Antonio Caso y José Vasconcelos desplazan, desde distintos planteamientos, el centro del problema hacia la persona y su proyecto cultural \citep{casoExistencia1916,vasconcelosRazaCosmica1925}. Frente a una visión exclusivamente cientificista o instrumental, recuperan dimensiones que no se agotan en el cálculo: desinterés, cultura, espiritualidad, identidad y proyecto colectivo. En consecuencia, esta aportación coincide con una idea básica de la ética filosófica: la acción humana requiere razones, pero también una comprensión del sentido de la vida y del lugar del otro \citep{huertaEticaConClasicos2000,singerCompendioEtica1995}.

Samuel Ramos introduce una mirada más introspectiva: la \textbf{ética de la identidad} también exige reconocer los rasgos culturales que condicionan la conducta \citep{ramosPerfil1934}. Su aporte no consiste en justificar defectos nacionales, sino en convertir la identidad en un problema filosófico. Para la práctica jurídica, esto es relevante porque las normas actúan sobre personas situadas en contextos históricos, sociales y culturales; ignorar ese contexto puede volver formalista cualquier decisión.

Con Sanchez Vazquez y Villoro, la reflexion etica mexicana adquiere un tono mas critico y politico. Sanchez Vazquez entiende la moral como comportamiento social y subraya la relacion entre libertad, praxis y transformacion \citep{sanchezVazquezEtica1969}. Villoro, por su parte, permite pensar la relacion entre poder y valor: una institucion puede ser legal y, aun asi, requerir examen etico si reproduce dominacion o excluye voces. Esta idea resulta especialmente importante para el Derecho, porque obliga a preguntar no solo si una decision esta permitida, sino si es justificable frente a la justicia y la dignidad.

Finalmente, Juliana Gonzalez Valenzuela representa una apertura hacia problemas contemporaneos de libertad, dignidad, cuerpo, vida y bioetica. Su importancia esta en recordar que la etica actual debe dialogar con avances cientificos, dilemas medicos, derechos humanos y responsabilidad profesional. En ese punto, la etica deja de ser una materia teorica para convertirse en criterio de decision frente a casos donde la vida humana, la autonomia y el cuidado estan en juego.

\section{Conclusiones}

La reflexion etica en Mexico ha evolucionado de una preocupacion educativa y civica hacia una discusion mas amplia sobre persona, identidad, justicia, poder, profesion y vida. El cambio no elimina las etapas anteriores; las acumula. Primero se pregunta como formar ciudadanos, despues como comprender a la persona y la cultura, y finalmente como actuar con responsabilidad ante instituciones, desigualdades y dilemas contemporaneos.

Entre los autores seleccionados hay coincidencias claras: todos entienden la etica como algo relacionado con la formacion humana y con la vida social. Sin embargo, difieren en el punto de partida. Barreda y Sierra privilegian la educacion; Caso y Vasconcelos, el humanismo; Ramos, la identidad; Sanchez Vazquez, la praxis social; Villoro, la justicia politica; Gonzalez, la dignidad y la responsabilidad ante la vida.

Considero que la relevancia profesional de este recorrido consiste en aprender a no separar Derecho y ética: la decisión jurídica debe examinarse no solo por su validez formal, sino también por sus efectos sobre las personas, la distribución del poder y la justicia \\citep{villoroPoderValor1997}. Como futuro abogado, sostengo que no basta con conocer normas; la responsabilidad profesional exige reconocer que toda decisión afecta derechos, distribuye cargas y expresa valores \\citep{ronquilloarmasEticaGeneralProfesional2018}. Por ello, la línea de tiempo permite concluir que la ética mexicana no ofrece una respuesta única, sino una tradición de preguntas que debe orientar la práctica jurídica: qué tipo de persona forma la educación, qué justicia sostiene la institución, qué dignidad protege la norma y qué responsabilidad asume quien decide \\citep{gonzalezEthosDestino1996,sanchezVazquezEtica1969}.\footnote{En la elaboración de este trabajo se utilizaron herramientas de inteligencia artificial como apoyo para la organización y revisión del texto, conforme a los Lineamientos para el uso ético y responsable de la Inteligencia Artificial en el ámbito académico de la UnADM. El análisis, la selección de fuentes y la postura sostenida son de mi autoría.}

\bibliography{etica-y-moral-juridica}

\end{document}


